% ============================================================================
%  INTRODUCTION GÉNÉRALE
% ============================================================================
\chapter*{Introduction générale}
\addcontentsline{toc}{chapter}{Introduction générale}
\markboth{Introduction générale}{Introduction générale}

Toute la théorie des communications construite depuis 1948 repose sur une idée de Shannon : transmettre, c'est
reproduire au récepteur, aussi exactement que possible, la suite de symboles choisie à l'émetteur, sans se
préoccuper de ce que ces symboles veulent dire~\cite{shannon1948}. Weaver distinguait pourtant déjà trois
niveaux de problème — le niveau technique (les symboles arrivent-ils ?), le niveau sémantique (le sens
arrive-t-il ?) et le niveau de l'efficacité (le destinataire agit-il correctement ?) — et constatait que seul le
premier avait été résolu~\cite{weaver1949}.

Ce découpage a produit l'architecture que nous utilisons tous les jours : on compresse la source d'un côté
(JPEG, H.264), on protège les bits de l'autre (code correcteur, modulation), et le théorème de séparation
garantit qu'aucune optimalité n'est perdue à les traiter séparément. La garantie est pourtant
\emph{asymptotique} : elle suppose des blocs infiniment longs et un canal dont on connaît la statistique.
Les images sont courtes, et le rapport signal sur bruit réel n'est jamais exactement celui qui a été prévu.

\paragraph{Problème.} Deux conséquences se manifestent. D'abord l'\textbf{effet falaise} : au-dessus du SNR de
conception la chaîne sature sans exploiter la marge disponible, et en dessous le décodeur décroche et l'image
est perdue d'un seul coup — pas dégradée, perdue. Ensuite le \textbf{mauvais critère} : la chaîne optimise la
fidélité pixel, alors que le destinataire est de plus en plus souvent une machine qui cherche seulement à
reconnaître le contenu. La question posée par ce projet est donc la suivante : \emph{à budget radio égal,
vaut-il mieux transmettre des pixels fidèles ou transmettre le sens ?}

\paragraph{Démarche et contributions.} Nous avons construit et mesuré trois systèmes complets dans des
conditions strictement identiques — même image, même rapport de bande, même canal, mêmes métriques. Le premier
est la chaîne séparée classique, simulée bit à bit et non approchée. Le deuxième est un codage conjoint
source-canal profond entraîné sur une perte pixel, qui reprend l'architecture de référence de la
littérature~\cite{bourtsoulatze2019}. Le troisième, notre proposition, est strictement le même réseau et ne
diffère que par la fonction de perte, à laquelle nous ajoutons un terme de caractéristiques et un terme de
tâche calculés à travers un classifieur gelé. Comme seule la perte change, tout écart de résultat lui est
imputable. Deux extensions complètent l'étude : une transmission orientée tâche et un renouvellement
conditionnel des latents en vidéo. L'ensemble est accompagné d'un simulateur interactif qui déroule une
transmission réelle étape par étape.

\paragraph{Organisation du rapport.} Le chapitre~\ref{chap:contexte} présente les fondements théoriques et
l'état de l'art. Le chapitre~\ref{chap:architecture} formule le problème, les exigences et l'architecture
retenue. Le chapitre~\ref{chap:separee} décrit le canal simulé et la chaîne séparée de référence. Le
chapitre~\ref{chap:jscc} détaille les systèmes appris et la perte sémantique. Le chapitre~\ref{chap:protocole}
expose le protocole d'évaluation et les métriques, en insistant sur l'indépendance des juges. Le
chapitre~\ref{chap:resultats} présente et analyse les résultats, y compris ceux qui nous sont défavorables. Le
chapitre~\ref{chap:extensions} traite les deux extensions, et le chapitre~\ref{chap:simulateur} l'outil
interactif développé pour explorer et vérifier la chaîne.
